%------------------------------------------------------------------
% Questão 4.21 – Intervalos de números com algarismos significativos
%------------------------------------------------------------------

\textbf{Ideia:} O intervalo vem da incerteza no último algarismo significativo. A incerteza é $\pm\frac{1}{2}$ da unidade daquele dígito.

%------------------------------------------------------------------
\textbf{(a) $7{,}7$}

\begin{itemize}
\item Algarismos significativos: 2
\item Último dígito: casa dos décimos (unidade = 0,1)
\item Incerteza: $\pm 0{,}05$
\end{itemize}

\[
\boxed{7{,}65 \leq 7{,}7 < 7{,}75}
\]

%------------------------------------------------------------------
\textbf{(b) $7{,}70$}

\begin{itemize}
\item Algarismos significativos: 3 (o zero final conta!)
\item Último dígito: casa dos centésimos (unidade = 0,01)
\item Incerteza: $\pm 0{,}005$
\end{itemize}

\[
\boxed{7{,}695 \leq 7{,}70 < 7{,}705}
\]

%------------------------------------------------------------------
\textbf{(c) $1200$}

\begin{itemize}
\item Algarismos significativos: 2 (assume-se 1 e 2 apenas)
\item Último dígito: casa das centenas (unidade = 100)
\item Incerteza: $\pm 50$
\end{itemize}

\[
\boxed{1150 \leq 1200 < 1250}
\]

\textit{Nota:} Para evitar ambiguidade, escreva $1{,}2 \times 10^3$.

%------------------------------------------------------------------
\textbf{(d) $1{,}200 \times 10^{-3}$}

\begin{itemize}
\item Algarismos significativos: 4
\item Último dígito: $10^{-6}$ (milionésimos)
\item Incerteza: $\pm 0{,}0000005 = \pm 5 \times 10^{-7}$
\end{itemize}

\[
\boxed{1{,}1995 \times 10^{-3} \leq 1{,}200 \times 10^{-3} < 1{,}2005 \times 10^{-3}}
\]

Ou em decimal: $0{,}0011995 \leq 0{,}001200 < 0{,}0012005$